% ============================================================
% setup/commands.tex —— 自定义命令 / 定理环境 / 算法汉化
% ============================================================

% ---------- 占位标记 ----------
% 正文中所有待写内容统一用 \placeholder{} 标出，终稿前逐项替换并删除。
% 占位文本为灰色斜体，便于一眼识别未完成部分。
\newcommand{\placeholder}[1]{%
  \textcolor{gray}{\textit{【占位：#1】}}}

% 写作要点提示（仅注释用，不输出）
% \writingnote{...} —— 以注释形式写在各节文件里更直观，此命令备用
\newcommand{\writingnote}[1]{}

% ---------- 首页页脚 ----------
% 无编号脚注：用于首页页脚的基金项目 / 通信作者（两刊投稿须知要求）
\newcommand{\blfootnote}[1]{%
  \begingroup
  \renewcommand{\thefootnote}{}%
  \footnote{#1}%
  \addtocounter{footnote}{-1}%
  \endgroup}

% ---------- 常用缩写/排版命令 ----------
\newcommand{\ie}{\textit{i.e.}}
\newcommand{\eg}{\textit{e.g.}}
\newcommand{\etal}{\textit{et al.}}
\newcommand{\R}{\mathbb{R}}          % 实数域
\newcommand{\abs}[1]{\left|#1\right|}
\newcommand{\norm}[1]{\left\|#1\right\|}
\newcommand{\transpose}{\mathrm{T}}      % 转置符号：正体 T（两刊均要求）

% ---------- 定理类环境（汉化） ----------
\theoremstyle{plain}
\newtheorem{definition}{定义}
\newtheorem{theorem}{定理}
\newtheorem{lemma}{引理}
\newtheorem{proposition}{命题}
\newtheorem{assumption}{假设}
\theoremstyle{remark}
\newtheorem{remark}{注}

% ---------- 算法浮动体汉化 ----------
\floatname{algorithm}{算法}
\algrenewcommand{\algorithmicrequire}{\textbf{输入：}}
\algrenewcommand{\algorithmicensure}{\textbf{输出：}}
